\documentclass{situationist-preprint}

\title{The Limits of Introspection}
\msid{MND-20-a}
\msfamily{MND-20}
\msdiscipline{Philosophy of Mind}
\mstype{Principal manuscript}
\mskeywords{introspection; self-knowledge; confabulation; phenomenal states; reliability; Nisbett and Wilson; Schwitzgebel}

\begin{document}
\maketitle

\begin{abstract}
Two literatures pull in opposite directions. Philosophers since Descartes have held that one knows one's own current experiences with a security unavailable for anything else; psychologists since Nisbett and Wilson have shown that people's reports of why they feel, choose and act as they do are often confabulated. Schwitzgebel has argued that the unreliability extends to current experience itself. We argue that introspection is reliable about current phenomenal states and unreliable about their causes. The difference is one of method. A judgment about the character of a current experience can be formed by a process in which the experience itself fixes the content of the judgment; a judgment about what caused the experience cannot, and must be reached by causal inference from a theory, which is the method an observer would use and is no more reliable in the first person than in the third. We show that the evidence of confabulation concerns causes, that Schwitzgebel's cases concern either non-current states, fine-grained structure beyond attention, or classification under theory-laden categories, and that a restricted claim of reliability survives them.
\end{abstract}

\section{The problem}

Descartes held that, whatever else might be doubted, it cannot be doubted that I seem to see light and hear noise (\emph{Meditations} II). \citet{Armstrong1968} gave the claim a naturalistic form, on which introspection is an inner sense, a causal mechanism for detecting one's own mental states, fallible in principle but reliable in fact. \citet{Gertler2011} surveys the accounts on which introspective knowledge is distinctively secure.

Against this, \citet{Nisbett1977} reviewed experiments in which subjects' explanations of their own choices and evaluations were systematically wrong. Subjects asked to choose among identical stockings preferred the rightmost pair and gave reasons about quality; subjects whose judgments had been influenced by a manipulation denied the influence and cited factors that had none. Strikingly, the subjects' explanations matched those that observers, told only of the situation, predicted. Later work extended the pattern. In choice blindness, subjects who were covertly handed the face they had not chosen defended the choice with reasons that referred to features of the wrong face \citep{Johansson2005}. In split-brain patients, the verbal hemisphere confidently explains actions initiated by the other on the basis of information it does not have \citep{Gazzaniga1978}. Misattribution studies show subjects ascribing arousal to the wrong source \citep{Schachter1962,Dutton1974}. \citet{Wilson2002} draws the moral that much of the mind is opaque to its owner, and \citet{Carruthers2011} argues that knowledge of one's own propositional attitudes is interpretive.

\citet{Schwitzgebel2008,Schwitzgebel2011} goes further: introspection of current conscious experience, he argues, is unreliable too. People are uncertain and in disagreement about whether they experience peripheral vision as clear, whether thought has a phenomenology, whether they dream in color, and what their emotions feel like. \citet{Hurlburt2007} stage a debate between Schwitzgebel and a defender of trained introspective sampling.

Our thesis is that the two kinds of evidence point to a single boundary. Introspection is reliable about current phenomenal states and unreliable about their causes.

\section{Framework}

\begin{definition}[Introspective judgment]\label{d:intro}
An \emph{introspective judgment} is a judgment about one's own mental life formed without reliance on observation of one's behavior or on testimony.
\end{definition}

\begin{definition}[Reliability]\label{d:rel}
A method is \emph{reliable} about a class of questions when, under normal conditions, most of the judgments it yields in answer to questions of that class are true.
\end{definition}

\begin{definition}[Current phenomenal state; cause]\label{d:classes}
A \emph{current phenomenal state} is a phenomenal property of an experience one is undergoing at the time of judgment, at the grain at which it is attended. A \emph{cause} of a mental state is any event or condition on which its occurrence or character counterfactually depends.
\end{definition}

The restriction to attended grain matters. MND-24 holds that phenomenal consciousness extends beyond attention; it follows that there are current phenomenal states about which no introspective judgment is made, and others, at the periphery, about which only coarse judgments can be made. The claim of reliability is not a claim about these.

\begin{principle}[Embedding]\label{p:embed}
A judgment about a current phenomenal state can be formed by a process in which the state itself constitutes part of the judgment's content, as when one judges \enquote{this is what I am now feeling}, with \enquote{this} fixed by the attended experience \citep{Chalmers2003}.
\end{principle}

\begin{principle}[Causal inference]\label{p:cause}
Knowledge of a causal relation between distinct events requires evidence of counterfactual dependence, obtained by observation of covariation or by a theory supported by such observation.
\end{principle}

Principle~\ref{p:embed} is the thesis, defended by Chalmers, that direct phenomenal concepts are formed by demonstrating an experience and that the judgments they figure in are partly constituted by the experience. Principle~\ref{p:cause} is a commonplace about causal knowledge: one event's being experienced does not show what produced it.

\section{The argument}

\begin{lemma}\label{l:states}
Introspective judgments about current phenomenal states, formed by the embedding method, are reliable.
\end{lemma}

\begin{proof}
By Principle~\ref{p:embed}, in a judgment formed by embedding the experience fixes what is judged: the content is that the attended state has the character it has, presented by the state itself. Such a judgment can go wrong in two ways only: if the demonstration fails to pick out a state, or if the subject goes on to classify the demonstrated state under a general concept, such as \enquote{pain} or \enquote{red}, and the classification is mistaken. The first failure requires abnormal conditions, since under normal conditions an attended experience is there to be demonstrated. The second concerns the application of general concepts to a given instance, which is reliable for the coarse categories that subjects are trained to apply and that are fixed by the experiences themselves, such as the difference between pain and pressure or between red and green. Hence most such judgments, under normal conditions, are true.
\end{proof}

\begin{lemma}\label{l:causes}
Introspective judgments about the causes of mental states are not reliable.
\end{lemma}

\begin{proof}
The cause of an experience is, in general, an event distinct from it. By Principle~\ref{p:cause}, a judgment that one caused the other requires evidence of counterfactual dependence. The experience does not carry such evidence: embedding the experience in the judgment fixes what it is like, not what it depends on. Introspection, by Definition~\ref{d:intro}, adds no observation of covariation. The introspector must therefore draw on a theory of what causes what in minds like hers, the same folk theory an observer uses, as Nisbett and Wilson's match between subjects' and observers' explanations shows. Where the theory is wrong, as it is for position effects, priming and misattributed arousal, introspective causal judgments are systematically wrong. Since such factors are pervasive, introspective judgments of cause are not reliable across the class.
\end{proof}

\begin{theorem}\label{t:main}
Introspection is reliable about current phenomenal states and unreliable about their causes.
\end{theorem}

\begin{proof}
From Lemmas~\ref{l:states} and \ref{l:causes}.
\end{proof}

\section{Schwitzgebel's cases}

Theorem~\ref{t:main} must withstand Schwitzgebel's argument that introspection fails for current experience itself. His cases fall into three groups, none of which is within the scope of Lemma~\ref{l:states}.

\paragraph{Non-current states.} Whether one dreams in color, and what one's emotions have generally felt like, are questions about past experience answered from memory and from theory. MND-25 holds that dream reports are partly reconstructed on waking. Unreliability here is a fact about memory and reconstruction, not about introspection of current experience.

\paragraph{Fine structure beyond attention.} How clear the visual field is a few degrees from fixation is a question about states at the periphery of attention, or about the detailed structure of the whole field. Definition~\ref{d:classes} restricts the claim to the grain at which a state is attended, and MND-24 explains why the periphery escapes it. When subjects attend to a point in the periphery, they report blur, and the reports are accurate.

\paragraph{Theory-laden classification.} Whether thought has a distinctive phenomenology, or whether an emotion is constituted by bodily feelings, are questions about how to classify an experience under categories whose application depends on a theory of the mind. Disagreement about them reflects disagreement in theory, as the long dispute over cognitive phenomenology illustrates; MND-19 settles that dispute in favor of a phenomenology of thinking by argument, not by an appeal to introspective reports alone. The embedding judgment that this is how one is now thinking is not in dispute.

\section{Objections and replies}

\begin{objection}[Choice blindness]\label{o:choice}
Subjects in choice-blindness studies misreport not only why they chose, but what they chose and what they saw. That is an error about current states.
\end{objection}

\begin{reply}
The subjects misremember which face they preferred a few seconds earlier and confabulate reasons for the choice. Both are errors about the past and about causes. Nothing in the studies suggests that subjects misdescribe the face they are now looking at, as it now looks to them.
\end{reply}

\begin{objection}[Pain and expectation]\label{o:pain}
Expectation alters how intense a pain is reported to be. If the report shifts while the stimulus is constant, introspection of pain is unreliable.
\end{objection}

\begin{reply}
The objection assumes that the stimulus fixes the pain. Expectation alters the experience itself, as placebo analgesia shows, and the report tracks the altered experience. What subjects cannot tell by introspection is that expectation was the cause of the change, which is what Theorem~\ref{t:main} predicts.
\end{reply}

\begin{objection}[Regress of classification]\label{o:class}
Every introspective report applies general concepts. If classification can be theory-laden, the line between Lemma~\ref{l:states} and the third group of cases is arbitrary.
\end{objection}

\begin{reply}
The line is drawn by how the concepts are applied. Concepts such as \enquote{pain} and \enquote{red} are learned by being applied to experiences on the basis of how they feel, and their application to a given experience is checked by that feeling. Concepts such as \enquote{cognitive phenomenology} or \enquote{constituted by bodily feeling} are applied on the basis of theory. The distinction admits borderline cases, but it does not need sharp boundaries to support a claim about most judgments in each class.
\end{reply}

\begin{objection}[Attitudes]\label{o:attitudes}
Knowledge of one's own beliefs and desires is not knowledge of phenomenal states. If it is interpretive, as Carruthers holds, the theorem covers little of what matters in self-knowledge.
\end{objection}

\begin{reply}
The theorem is not an account of knowledge of attitudes, for which EPI-15 defends transparency: one knows whether one believes $p$ by considering whether $p$ \citep[see also][]{Byrne2018}. EPI-15 also holds that the transparency method fails for attitudes not answerable to reasons, such as a recalcitrant fear; for these, introspection of how one currently feels is one of the routes that remain, and Theorem~\ref{t:main} says that it is reliable about how one feels, though not about why.
\end{reply}

\section{What the result does not establish}

The result does not establish that introspection is infallible about current experience, or that it is reliable under abnormal conditions, such as intoxication, hypnosis or some neurological disorders. It does not cover states outside attention or the fine structure of the experiential field. It does not show that causal self-knowledge is impossible: one can learn the causes of one's states by observation and experiment, as anyone else can, and the claim is only that introspection does not supply it.

\section{Conclusion}

Introspection is reliable about current phenomenal states and unreliable about their causes. The evidence of confabulation shows that people answer questions about the causes of their states with a folk theory, as an observer would. Introspection of the current states themselves is different, because the states fix the content of the judgments that report them. Schwitzgebel's cases fall outside this domain: they concern memory, the unattended periphery or theory-laden classification.

\begin{thebibliography}{99}\small
\bibitem[Armstrong(1968)]{Armstrong1968} Armstrong, D.~M. (1968). \emph{A Materialist Theory of the Mind}. Routledge and Kegan Paul.
\bibitem[Byrne(2018)]{Byrne2018} Byrne, A. (2018). \emph{Transparency and Self-Knowledge}. Oxford University Press.
\bibitem[Carruthers(2011)]{Carruthers2011} Carruthers, P. (2011). \emph{The Opacity of Mind: An Integrative Theory of Self-Knowledge}. Oxford University Press.
\bibitem[Chalmers(2003)]{Chalmers2003} Chalmers, D.~J. (2003). The content and epistemology of phenomenal belief. In Q. Smith and A. Jokic (eds.), \emph{Consciousness: New Philosophical Perspectives}, 220--272. Oxford University Press.
\bibitem[Dutton and Aron(1974)]{Dutton1974} Dutton, D.~G. and Aron, A.~P. (1974). Some evidence for heightened sexual attraction under conditions of high anxiety. \emph{Journal of Personality and Social Psychology}, 30(4), 510--517.
\bibitem[Gazzaniga and LeDoux(1978)]{Gazzaniga1978} Gazzaniga, M.~S. and LeDoux, J.~E. (1978). \emph{The Integrated Mind}. Plenum Press.
\bibitem[Gertler(2011)]{Gertler2011} Gertler, B. (2011). \emph{Self-Knowledge}. Routledge.
\bibitem[Hurlburt and Schwitzgebel(2007)]{Hurlburt2007} Hurlburt, R.~T. and Schwitzgebel, E. (2007). \emph{Describing Inner Experience? Proponent Meets Skeptic}. MIT Press.
\bibitem[Johansson et~al.(2005)]{Johansson2005} Johansson, P., Hall, L., Sikstr\"om, S. and Olsson, A. (2005). Failure to detect mismatches between intention and outcome in a simple decision task. \emph{Science}, 310(5745), 116--119.
\bibitem[Nisbett and Wilson(1977)]{Nisbett1977} Nisbett, R.~E. and Wilson, T.~D. (1977). Telling more than we can know: Verbal reports on mental processes. \emph{Psychological Review}, 84(3), 231--259.
\bibitem[Schachter and Singer(1962)]{Schachter1962} Schachter, S. and Singer, J.~E. (1962). Cognitive, social, and physiological determinants of emotional state. \emph{Psychological Review}, 69(5), 379--399.
\bibitem[Schwitzgebel(2008)]{Schwitzgebel2008} Schwitzgebel, E. (2008). The unreliability of naive introspection. \emph{The Philosophical Review}, 117(2), 245--273.
\bibitem[Schwitzgebel(2011)]{Schwitzgebel2011} Schwitzgebel, E. (2011). \emph{Perplexities of Consciousness}. MIT Press.
\bibitem[Wilson(2002)]{Wilson2002} Wilson, T.~D. (2002). \emph{Strangers to Ourselves: Discovering the Adaptive Unconscious}. Harvard University Press.
\end{thebibliography}

\end{document}
